\documentclass{article}
\usepackage[utf8]{inputenc}
\usepackage[T1]{fontenc}
\usepackage{geometry}
\usepackage{graphicx}
\usepackage{float}
\usepackage{hyperref}
\usepackage[english,indonesian]{babel}

\title{Optimasi Pemrosesan Template LaTeX Menggunakan Pendekatan Berbasis AI}
\author{Academic Support Agent}
\date{\today}

\begin{document}
\maketitle

\begin{abstract}

\end{abstract}

\chapter{PENDAHULUAN}


\chapter{TINJAUAN PUSTAKA}


% Original content (commented out):
% 
\documentclass[12pt]{article}
 
\usepackage[margin=1in]{geometry} 
\usepackage{amsmath,amsthm,amssymb,scrextend}
\usepackage{fancyhdr}
\pagestyle{fancy}
\DeclareMathOperator{\rng}{Rng}
\DeclareMathOperator{\dom}{Dom}
\newcommand{\R}{\mathbb R}
\newcommand{\cont}{\subseteq}
\newcommand{\N}{\mathbb N}
\newcommand{\Z}{\mathbb Z}
\usepackage{tikz}
\usepackage{pgfplots}
\usepackage{amsmath}
\usepackage[mathscr]{euscript}
\let\euscr\mathscr \let\mathscr\relax% just so we can load this and rsfs
\usepackage[scr]{rsfso}
\usepackage{amsthm}
\usepackage{amssymb}
\usepackage{multicol}
\usepackage[colorlinks=true, pdfstartview=FitV, linkcolor=blue,
citecolor=blue, urlcolor=blue]{hyperref}

\DeclareMathOperator{\arcsec}{arcsec}
\DeclareMathOperator{\arccot}{arccot}
\DeclareMathOperator{\arccsc}{arccsc}
\newcommand{\ddx}{\frac{d}{dx}}
\newcommand{\dfdx}{\frac{df}{dx}}
\newcommand{\ddxp}[1]{\frac{d}{dx}\left( #1 \right)}
\newcommand{\dydx}{\frac{dy}{dx}}
\let\ds\displaystyle
\newcommand{\intx}[1]{\int #1 \, dx}
\newcommand{\intt}[1]{\int #1 \, dt}
\newcommand{\defint}[3]{\int_{#1}^{#2} #3 \, dx}
\newcommand{\imp}{\Rightarrow}
\newcommand{\un}{\cup}
\newcommand{\inter}{\cap}
\newcommand{\ps}{\mathscr{P}}
\newcommand{\set}[1]{\left\{ #1 \right\}}
\newtheorem*{sol}{Solution}
\newtheorem*{claim}{Claim}
\newtheorem{problem}{Problem}
\begin{document}
 
% EVERYTHING ABOVE THIS LINE IS JUST PREABLE, NO NEED TO MESS WITH IT.__________________________________________________________________________________________
%

\lhead{Math 108}
\chead{Your Name}
\rhead{\today}
 
% Just put your proofs in between the \begin{proof} and the \end{proof} statements!

\section*{Homework 7}
\begin{problem} Let $S=\set{(x,y)\in \R \times \R: y=\sqrt{4-2x}}$ prove that: 
\begin{enumerate}
\item $0\in \rng(S)$
\begin{proof} 

\end{proof}
\item $3\not\in \dom(S)$ 
\begin{proof} 

\end{proof}
\end{enumerate}

\end{problem}
 

\begin{problem} 
Let $R=\set{(x,y)\in \mathbb R \times \mathbb R: x=|y|} $, prove that:
\begin{enumerate}
\item $\dom(R)=[0,\infty)$ 
\begin{proof} 

\end{proof}
\item $\rng(R)=\R$ 
\begin{proof} 

\end{proof}
\end{enumerate}
\end{problem}

\begin{problem} Prove that $\dom(S \circ R ) \cont \dom(R)$
\end{problem}
\begin{proof}

\end{proof}
\begin{problem}
One of these statements is true, the other is false. Prove the statement that is true and give a counter example to show the other statement is false:
\begin{enumerate}
\item $\rng(S) \cont \rng(S \circ R)$
\begin{proof}

\end{proof}
\item $\rng(S \circ R) \cont \rng(S)$
\begin{proof}

\end{proof}
\end{enumerate}
\end{problem}

\begin{problem}
Prove $(R^{-1})^{-1} = R$
\end{problem}
\begin{proof}

\end{proof}
% THE DOCUMENT IS ESSENTIALLY DONE AT THIS POINT, NO NEED TO EDIT ANYTHING BELOW THIS______________________________________________________________________________________________
 
\end{document}

\end{document}